\chapter{Conclusione}

\section{Impressioni sul progetto}
Ho sempre voluto scrivere un progetto del genere perche' ero interessato a capire veramente il
funzionamento degli applicativi che utilizziamo tutti i giorni come WhatsApp che trovo molto
interessanti. Mi sono divertito a cercare gli algoritmi e imparare le implementazioni. Sicuramente
la parte piu' difficile e' stato lo strutturare bene il codice e cercare di renderlo pulito. In
generale sono soddisfatto del prodotto ottenuto ma sono consapevole che non sia un prodotto finito
ma che richieda ancora lavoro per essere considerato utilizzabile. E' stato un bel percorso molto
interessante che mi ha permesso di imparare diverse cose tipo il discorso sugli hash delle password
a cui non avevo pensato.

\section{Sviluppi futuri}
Ci sono moltissime cose da migliorare se volessi rendere questa applicazione veramente pubblica e
utilizzabile.

\subsection{Algoritmi di cifratura}
Sicuramente la prima cosa da fare e' rendere veramente sicura la comunicazione implementando AES per
la cifratura simmetrica e probabilmente Diffie Hellman sulle curve ellittiche per lo scambio delle
chiavi. Inoltre si dovrebbe mettere la chiave RSA a 4096 bit.
Un'altra cosa interessante da aggiungere sarebbe di cambiare le chiavi di sessione ogni N messaggi.

\subsection{Comunicazione tra client}
Al momento la comunicazione tra client non garantisce la forward secrecy perche' utilizzo RSA per
cifrare la chiave. Potrei utilizzare la stessa implementazione di client e server ma questo
implicherebbe che per inviare un messaggio a un nuovo utente egli debba essere online visto che
Diffie Hellman richiede cooperazione. Una cosa realistica che si potrebbe implementare in futuro e'
un sistema per cui ogni client carica sul server un certo numero di chiavi precalcolate e quando
qualcuno vuole creare una comunicazione prende una di quelle chiavi usa e getta e si crea la chiave
privata senza avere il bisogno che il ricevente sia online.

\subsection{Linguaggio}
Per ottenere la sicurezza dell'applicativo non e' necessario cambiare linguaggio ma sicuramente
utilizzerei delle librerie gia' fatte che eseguono le istruzioni in C, ad esempio pycryptodome.
Questo perche' mi rendo conto che il mio codice e' lentissimo a generare i numeri primi grandi e
purtroppo questo accade principalmente perche' ho usato Python puro senza chiamate a C.

\subsection{UI}
Sicuramente un miglioramento alla UI e' necessario in quanto non le ho dato alcun peso durante il
progetto, ho semplicemente usato la libreria tkinter per rendere l'implementazione piu' rapida
mettendo qualche bottone e qualche textbox. Un ottimo update sarebbe di usare html, js e css per la
UI.

\subsection{Funzionalita'}
Una volta creata una buona base si puo' pensare a integrare diverse funzionalita' come lo scambio di
file e la creazione di gruppi.
